% !TEX TS-program = XeLaTeX
% use the following command:
% all document files must be coded in UTF-8
\documentclass[portuguese]{textolivre}
% build HTML with: make4ht -e build.lua -c textolivre.cfg -x -u article "fn-in,svg,pic-align"

\journalname{Texto Livre}
\thevolume{19}
%\thenumber{1} % old template
\theyear{2026}
\receiveddate{\DTMdisplaydate{2025}{11}{7}{-1}} % YYYY MM DD
\accepteddate{\DTMdisplaydate{2026}{1}{7}{-1}}
\publisheddate{\DTMdisplaydate{2026}{5}{13}{-1}}
\corrauthor{Ana Elisa Ribeiro}
\articledoi{10.1590/1983-3652.2026.62700}
%\articleid{NNNN} % if the article ID is not the last 5 numbers of its DOI, provide it using \articleid{} commmand 
% list of available sesscions in the journal: articles, dossier, reports, essays, reviews, interviews, editorial
\articlesessionname{dossier}
\runningauthor{Ribeiro} 
%\editorname{Leonardo Araújo} % old template
\sectioneditorname{Francis Arthuso Paiva~\orcid{0000-0002-9083-3342}}
\layouteditorname{Saula Cecília~\orcid{0009-0006-3069-8480}}

\title{Definições atualizadas e uso com moderação: um ensaio sobre inteligência artificial na Educação Básica}
\othertitle{Updated definitions and use with moderation: an essay on artificial intelligence in primary and secondary education}
% if there is a third language title, add here:
%\othertitle{Artikelvorlage zur Einreichung beim Texto Livre Journal}

\author[1]{Ana Elisa Ribeiro~\orcid{0000-0002-4422-7480}\thanks{Email: \href{mailto: anadigital@gmail.com}{anadigital@gmail.com}}}
\affil[1]{Centro Federal de Educação Tecnológica de Minas Gerais (CEFET-MG), Departamento de Linguagem e Tecnologia, Belo Horizonte, MG, Brasil.}

\addbibresource{article.bib}
% use biber instead of bibtex
% $ biber article

% used to create dummy text for the template file
\definecolor{dark-gray}{gray}{0.35} % color used to display dummy texts
\usepackage{lipsum}
\SetLipsumParListSurrounders{\colorlet{oldcolor}{.}\color{dark-gray}}{\color{oldcolor}}

% used here only to provide the XeLaTeX and BibTeX logos
\usepackage{hologo}

% if you use multirows in a table, include the multirow package
\usepackage{multirow}

% provides sidewaysfigure environment
\usepackage{rotating}

% CUSTOM EPIGRAPH - BEGIN 
%%% https://tex.stackexchange.com/questions/193178/specific-epigraph-style
\usepackage{epigraph}
\renewcommand\textflush{flushright}
\makeatletter
\newlength\epitextskip
\pretocmd{\@epitext}{\em}{}{}
\apptocmd{\@epitext}{\em}{}{}
\patchcmd{\epigraph}{\@epitext{#1}\\}{\@epitext{#1}\\[\epitextskip]}{}{}
\makeatother
\setlength\epigraphrule{0pt}
\setlength\epitextskip{0.5ex}
\setlength\epigraphwidth{.7\textwidth}
% CUSTOM EPIGRAPH - END

% to use IPA symbols in unicode add
%\usepackage{fontspec}
%\newfontfamily\ipafont{CMU Serif}
%\newcommand{\ipa}[1]{{\ipafont #1}}
% and in the text you may use the \ipa{...} command passing the symbols in unicode

% LANGUAGE - BEGIN
% ARABIC
% for languages that use special fonts, you must provide the typeface that will be used
% \setotherlanguage{arabic}
% \newfontfamily\arabicfont[Script=Arabic]{Amiri}
% \newfontfamily\arabicfontsf[Script=Arabic]{Amiri}
% \newfontfamily\arabicfonttt[Script=Arabic]{Amiri}
%
% in the article, to add arabic text use: \textlang{arabic}{ ... }
%
% RUSSIAN
% for russian text we also need to define fonts with support for Cyrillic script
% \usepackage{fontspec}
% \setotherlanguage{russian}
% \newfontfamily\cyrillicfont{Times New Roman}
% \newfontfamily\cyrillicfontsf{Times New Roman}[Script=Cyrillic]
% \newfontfamily\cyrillicfonttt{Times New Roman}[Script=Cyrillic]
%
% in the text use \begin{russian} ... \end{russian}
% LANGUAGE - END

% EMOJIS - BEGIN
% to use emoticons in your manuscript
% https://stackoverflow.com/questions/190145/how-to-insert-emoticons-in-latex/57076064
% using font Symbola, which has full support
% the font may be downloaded at:
% https://dn-works.com/ufas/
% add to preamble:
% \newfontfamily\Symbola{Symbola}
% in the text use:
% {\Symbola }
% EMOJIS - END

% LABEL REFERENCE TO DESCRIPTIVE LIST - BEGIN
% reference itens in a descriptive list using their labels instead of numbers
% insert the code below in the preambule:
%\makeatletter
%\let\orgdescriptionlabel\descriptionlabel
%\renewcommand*{\descriptionlabel}[1]{%
%  \let\orglabel\label
%  \let\label\@gobble
%  \phantomsection
%  \edef\@currentlabel{#1\unskip}%
%  \let\label\orglabel
%  \orgdescriptionlabel{#1}%
%}
%\makeatother
%
% in your document, use as illustraded here:
%\begin{description}
%  \item[first\label{itm1}] this is only an example;
%  % ...  add more items
%\end{description}
% LABEL REFERENCE TO DESCRIPTIVE LIST - END


% add line numbers for submission
%\usepackage{lineno}
%\linenumbers

\begin{document}
\maketitle

\begin{polyabstract}
\begin{abstract}
Este ensaio discute o surgimento das inteligências artificiais gerativas (IAG) e seu impacto sobre o ensino de Português nas escolas, em articulação com a BNCC e com a empiria da experiência em sala de aula na escola pública. Argumenta-se que a IAG nem sempre pode ser vista como ``ferramenta'' desejável na sala de aula de Português em razão do estágio da aprendizagem de produção de texto dos estudantes. O texto relata brevemente a história da IA, discute questões éticas e morais dos usos da IAG, notadamente na escola, e retoma questões como o erro, o desenvolvimento da autoria e o senso crítico na produção textual de crianças e jovens.

\keywords{Inteligência artificial gerativa\sep Produção de textos na escola\sep Tecnologias digitais e ensino\sep Ensino de Português}
\end{abstract}

\begin{english}
\begin{abstract}
This essay discusses the emergence of generative artificial intelligence and its impact on the teaching of Portuguese in schools, in conjunction with the BNCC (Brazilian National Curriculum Base) and the empirical experience in a public school classroom. It argues that AIGen may not always be seen as a desirable ``tool'' in the Portuguese classroom due to the stage of text production learning of students. The text briefly recounts the history of AI, discusses ethical and moral issues surrounding its use, particularly in schools, and revisits issues such as error, the development of authorship, and critical thinking in the text production of children and young people.

\keywords{Generative artificial intelligence\sep Text production in school\sep Digital technologies and education\sep Portuguese language teaching}
\end{abstract}
\end{english}
% if there is another abstract, insert it here using the same scheme
\end{polyabstract}

\section{Considerações iniciais: especular}\label{sec-intro}
Este é um texto especulativo. Ele se vale de resultados de pesquisa e de teorizações sobre as inteligências artificiais e o ensino de língua portuguesa, mas é, fundamentalmente, um exercício de imaginação ou de crítica em relação à adesão às tecnologias que vêm sendo chamadas de IA, e mais especificamente as IA Gerativas (IAG), isto é, capazes de gerar textos, imagens ou sons sob um comando conhecido como \textit{prompt}, palavra que se popularizou diante da divulgação ampla da existência e mesmo das interfaces de inteligências artificiais acessíveis a muitos usuários\footnote{A mais conhecida delas tornou-se o ChatGPT (sigla para Generative Pre-Trained Transformer), da OpenAI, desenvolvida em 2018-2019 e hoje bastante popular. Há outras empresas e chats semelhantes, mas, neste texto, evitarei marcas e preferirei tratá-las genericamente como IA ou IAG, elegendo a palavra ``gerativa'' para o G, e abarcando essas IA baseadas em grandes bancos de exemplos e modelos. Neste texto, também usarei como sinônimos ``produção de textos'' e ``redação'', a despeito de uma longa tradição de debates sobre isso no Brasil e das diferenças de abordagem conhecidas. ``Redação'' ainda é o nome que se dá ao texto solicitado aos estudantes pelo Exame Nacional do Ensino Médio (Enem) e é o nome oficial da disciplina que eu mesma ministro há muitos anos no ensino médio do CEFET-MG.}.

A rigor, talvez todo texto sobre IA e ensino seja especulativo. Diante da recente dispersão das IAG, tanto em relação à quantidade de empresas/programas disponíveis, quanto em relação à notícia da existência dessas tecnologias, é possível que textos muito definitivos sobre o assunto soem algo arriscados, até falaciosos, no que diz respeito a aspectos das IAG que ainda são pouco conhecidos e têm sido debatidos em espaços de educação realmente preocupados com a formação de crianças e jovens.

Neste ensaio, pretendo fazer alguns movimentos na direção de pensar sobre as inteligências artificiais, em especial as que geram textos verbais, na relação com o ensino de língua portuguesa, aqui entendido como o ensino de leitura e escrita, marcadamente na educação básica, considerando estágios da escolarização relacionados à idade dos alunos e das alunas, assim como os níveis de complexidade e sofisticação supostamente progressivos dos textos a serem lidos e escritos, principalmente estes, ao longo das etapas escolares, isto é, o ensino fundamental e o ensino médio.

\section{BNCC e tecnologias}
Para tratar do contexto escolar, é incontornável retomar a \textit{Base Nacional Comum Curricular} (BNCC) a fim de tratar dos direitos que as instituições educacionais deveriam assegurar aos jovens brasileiros. Para a área de Linguagens, de meu interesse direto aqui, é preciso destacar que o documento oficial preconiza, para o ensino fundamental, o desenvolvimento de linguagens visuais, sonoras, verbais e corporais, ``visando estabelecer um repertório diversificado sobre as práticas de linguagem e desenvolver o senso estético e a comunicação com o uso das tecnologias digitais'' \cite[p. 470]{brasil2018}. Para a etapa do ensino médio, a BNCC estabelece o investimento na

\begin{quote}
    ampliação da autonomia, do protagonismo e da autoria nas práticas de diferentes linguagens; na identificação e na crítica aos diferentes usos das linguagens, explicitando seu poder no estabelecimento de relações; na apreciação e na participação em diversas manifestações artísticas e culturais e no uso criativo das diversas mídias  \cite[p. 470]{brasil2018}.
\end{quote}

Do trecho citado, parece importante destacar aspectos como a autoria, a crítica e o uso criativo de mídias diversas, elementos que não devem ser esquecidos não apenas na construção desta argumentação, mas também, e principalmente, como orientadores das tarefas diárias de professoras e professores de português (leitura e redação). Segundo a BNCC, o ensino e a aprendizagem das Linguagens devem ``possibilitar uma participação mais plena dos jovens nas diferentes práticas sociais que envolvem o uso das linguagens'' \cite[p. 473]{brasil2018}. É de se destacar também que as tecnologias digitais (``mundo digital'', ``cultura digital'' etc.) estão muito presentes no documento, inclusive como prioridade, já que se entende que têm impacto no dia a dia das pessoas. Isso, é claro, já estava considerado antes de as IAG ocuparem os espaços de discussão e prática que ocupam hoje. Nos dias que correm, elas ganham o status de tema inescapável tanto nas salas de aula, quanto nos espaços de formação de professores. Pode-se dizer, sem tanto risco assim, que aprender sobre IAG é, hoje, importante para evitar o alheamento dos/as jovens, num mundo já profundamente afetado por essas tecnologias; melhorar os ``resultados sociais'' \cite{cazden1996}\footnote{A expressão, assim como a preocupação, está no conhecido manifesto da pedagogia dos multiletramentos assinado por dez docentes de países diversos, há trinta anos. O artigo foi publicado em 1996 e influenciou fortemente a discussão sobre letramentos e multiletramentos no Brasil. Uma tradução autorizada do documento foi publicada em acesso aberto pela Led, editora-laboratório do bacharelado em Letras do CEFET-MG (ver em \url{https://www.led.cefetmg.br/uma-pedagogia-dos-multiletramentos/}) \cite{cazden2021}. Uma discussão sobre a influência do manifesto no Brasil está em \textcite{ribeiro2020}.} do que ensinamos na escola (e nas aulas de línguas, especialmente); atentar para um cenário de imensas e mantidas desigualdades sociais, o que inclui crianças e jovens com ou sem acesso tanto às tecnologias, elas mesmas, quanto ao debate informado e crítico sobre elas; não perder de vista um cenário de desigualdade e assimetria de acessos e usos ainda maior.  

Segundo a \textit{Base}, a utilização de tecnologias digitais na escola ``não só possibilita maior apropriação técnica e crítica desses recursos, como também é determinante para uma aprendizagem significativa e autônoma pelos estudantes'' \cite[p. 478]{brasil2018}. A despeito desse direcionamento reiterado ao longo do documento, também é clara a orientação de que a ``cultura do impresso (ou da palavra escrita)'' \cite[p. 478]{brasil2018} deve continuar central na educação escolar. A manutenção desses dois horizontes de tecnologia da cultura escrita importa, mais do que nunca, se pensarmos em termos da tecnodiversidade \cite{hui2020} e de seu efeito resistente ao achatamento das opções e das possibilidades técnicas a que uma pessoa pode/deve ter acesso para estar em sociedade. Interessa também considerar que nossas definições de ``tecnologias digitais'' foram atualizadas (como nas mensagens de \textit{feedback} das máquinas...), e que as IA cabem agora sob esse vago rótulo, mesmo que ainda não tenhamos plena compreensão do que são e de seus efeitos e impactos sociais e educacionais, para mencionar apenas dois deles.

Este ponto interessa de perto: a centralidade da escrita. Afora certa confusão conceitual e técnica que o texto da BNCC permite flagrar entre impresso e escrito\footnote{ A cultura escrita abarca tecnologias impressas, manuscritas, digitais etc. Não é nem sinônima de tecnologias impressas, nem de digitais. A escrita continua central em nossa sociedade, por meio de técnicas e tecnologias, hoje mais diversas e mais imbricadas (diz-se que somos uma sociedade ``grafocêntrica''). Um livro impresso, por exemplo, provavelmente nasce da digitação em um editor de textos digital, mas ainda pode ter se originado em anotações manuscritas. Tudo isso concorre para a complexificação do que se chama de cultura escrita. Na BNCC, isso aparece bastante misturado.}, é preciso ter claro que escrever, inclusive de diversos modos e com variados recursos tecnológicos, segue sendo um conhecimento e uma atividade central para estudantes da educação básica. Visto de outra perspectiva, o ensinar a escrever se mantém como tarefa e missão incontornável do professor e da professora de língua, ao menos por enquanto. É prudente, embora incômodo, considerar que até mesmo essa missão possa deslizar – ou ser empurrada – para outros lugares. A questão é que o ensinar a ler e escrever não se dá mais, e faz tempo, a partir de um repertório apenas analógico de possibilidades (caneta, papel, brochura etc.), mas dentro de uma ``paisagem comunicacional'' \cite{kress2003} muito mais variada (teclados para digitar, telas de todos os tamanhos, por exemplo), que inclui, agora, as inteligências artificiais, se consideradas no rol das tecnologias digitais a serem compreendidas, consideradas, debatidas, experimentadas e, quem sabe, criticadas em sua relação com a produção escrita.

A BNCC toma como campos de atuação para a escola: as práticas de estudo e pesquisa (salientando a apreciação e a produção de textos analíticos e argumentativos); o campo jornalístico-midiático (com foco na circulação de textos das mídias impressa, audiovisual e digital, além da publicidade); a vida pública (que contempla textos normativos, legais e jurídicos, propositivos e reivindicatórios); e a criação artística (destacando a fruição e a sensibilidade estética). Em todos os casos, a leitura e a produção de textos são centrais, em suas manifestações e modos de circulação possíveis atualmente.

Nesse contexto de compreensão da BNCC e de sua tradução ao dia a dia do ensino de língua nas escolas regulares, consideradas as orientações explícitas e reiteradas sobre as interações entre leitura, escrita e tecnologias (analógicas e digitais), a IAG surge como um elemento de grande impacto, já que interfere diretamente no ensino da escrita e na produção de textos por estudantes de todas as etapas. Também tem seus efeitos para a leitura, já que pode gerar textos que criam uma ilusão de honestidade, até mesmo autoria ou humanidade, mas que precisam ser lidos e avaliados, inclusive analisados sobre sua condição de verdade, correção e confiabilidade. É, hoje, inescapável comentar sobre usos de IAG por estudantes (e não só eles...) na escola, sob práticas ainda pouco ou nada reguladas (em termos morais e éticos, por exemplo), com efeitos ainda desconhecidos, em termos de aprendizagem. Por hora, as IAG têm sido vistas como amigáveis, pelos mais otimistas, e como ameaças, pelos desconfiados. É, no entanto, possível afirmar que os efeitos de sua chegada às práticas sociais de comunicação (e à educação, que mais nos interessa aqui) ainda não puderam ser flagrados, analisados e compreendidos o bastante. Há certo alumbramento com a chegada das IAG e o que elas geram, o que me deixa em uma posição a um só tempo de conforto para especular e de risco para tecer afirmações.

\section{Inteligências artificiais gerativas}
Contar uma breve história das inteligências artificiais é importante para que se tenha consciência de sua relativa longa duração, assim como de sua evolução, em termos de interfaces, resultados e presença social. Segundo \textcite{coeckelbergh2023}, a disciplina conhecida como IA teve início nos anos 1950, depois de duas invenções: a do computador digital programável, na década anterior, e a dos estudos do que Norbert Wiener chamou de ``cibernética''. Nos anos 1950, houve outro marco importante: a publicação, na revista \textit{Mind}, do artigo de Alan Turing que apresentava o famoso teste que levava seu nome. A ideia era saber se as máquinas podiam pensar ou realizar tarefas abstratas. O termo ``inteligência artificial'', no entanto, só foi cunhado no final dos anos 1950, por John McCarthy. Naquele momento histórico, a diferença entre máquinas analógicas e máquinas digitais começava a despontar como um grande interesse de cientistas e professores.

Os tipos de IA existentes, até há pouco tempo, eram baseados em determinações introduzidas nos sistemas, e suas respostas eram mais ou menos previsíveis ou determinísticas (inclusive as desenvolvidas sob o paradigma do conexionismo, nos anos 1980)\footnote{Embora não seja minha intenção elencar os tipos de IA que existiram ou existem, pode ser importante mencionar o que se chamou de ``aprendizagem por máquina'' ou aplicativos hoje muito populares, tais como Waze, Alexa, entre outros, que se baseiam em IA. Também menciono, por sugestão de um(a) colega leitor(a) crítico(a) deste trabalho, a chatbot ELIZA, considerada a primeira robô de conversação, ainda nos anos 1960. Ela foi criada por Joseph Weizenbaum e simulava um processo de psicoterapia por meio de perguntas e respostas. Fazia isso tão bem e tão semelhantemente a um humano, que esse tipo de ocorrência ficou conhecido como ``efeito Eliza''.}. Somente nos anos 2010, as IA passaram a usar tecnologias transformativas que permitiram que algoritmos calculassem probabilidades a partir de imensos bancos e modelos linguísticos, o que finalmente resultou na chamada ``inteligência artificial gerativa'' ou ``generativa'' (em uma tradução mais próxima da palavra em inglês), isto é, uma IA que responde a um comando de texto (e em todos os casos, mesmo para respostas com imagens, a entrada se dá por meio de texto verbal) e entrega um resultado textual geralmente muito satisfatório em relação ao texto produzido por um humano. Não raro, melhor. Esse talvez seja o elemento que tem fascinado, se não deslumbrado, usuários do mundo inteiro, que se veem diante de uma interface simples, capaz de gerar um texto em segundos. É, no entanto, fundamental saber algo aparentemente simples: a IA não \textit{escreve}, isso tido como uma capacidade fundamentalmente humana; a IA \textit{gera} texto.

As aplicações das IA estão por todos os lugares. As IAG, isto é, geradoras de textos, imagens e até música, é que têm ocupado não apenas o noticiário, mas as salas de aula, em todos os níveis de ensino, com impactos ainda por medir ou conhecer. \textcite[p. 79]{coeckelbergh2023} acredita que haverá imenso ``impacto nas relações sociais, econômicas e ambientais'', causando alteração nas interações humanas. O autor afirma: ``Não é apenas sobre tecnologia, mas também sobre o que os seres humanos fazem com ela, como a usam, como a percebem e a experimentam, e como a incorporamos em ambientes sociotécnicos mais amplos''. O momento atual, em que a IAG causa grande perturbação, é também oportunidade de reflexão sobre uma educação em mídias e tecnologias sempre dinâmica, que insta professores e professoras ao debate sobre o ensino de leitura e escrita, assim como sobre a própria escrita e seus processos. \textcite[p. 79]{coeckelbergh2023} lembra que, historicamente, ``quando a tecnologia é nova, tende a ser altamente controversa, mas, uma vez que se torna incorporada na vida quotidiana, a euforia e a controvérsia diminuem significativamente''. O arrefecimento da excitação provocada pela chegada da IAG às palmas das mãos de muitas pessoas não parece estar próximo, por enquanto.

Para o caso das IA capazes de gerar textos a partir de comandos de texto \textit{(prompts)}, parece fundamental construir alguns conhecimentos sobre seu processo tecnológico, a fim de tecer relações produtivas e informadas sobre como elas funcionam, o que geram, como o fazem e em que medida revolvem a cena do ensino de escrita na escola básica. Um ponto relevante é saber que a IAG é um tipo de programa de computador que, a partir de entradas e regras nele introduzidas, gera saídas sem entendê-las, isto é, são máquinas sem intencionalidade e capacidade de juízo moral. Outro ponto importante é que há seres humanos implicados em todas as etapas dos processos da IA: estruturação de problemas, preparação de bancos de dados, treinamento de algoritmos etc. O conhecimento de especialistas humanos é indispensável para a operação das IAG, assim como os resultados por elas gerados não dispensam, de forma alguma, a avaliação e a verificação humanas, ao menos por enquanto.

\textcite{coeckelbergh2023} levanta uma questão de grande aplicação para a área educacional. O autor considera que alguns usuários de IA são mais vulneráveis do que outros em relação a vários pontos: dependência da tecnologia para tomada de decisões, aceitação acrítica de textos enviesados (replicando preconceitos, por exemplo), distanciamento em relação à autoria e à responsabilidade sobre resultados, entre outros. Os ``vulneráveis tecnológicos'' ficam à mercê dos resultados gerados pela IA e não têm a capacidade de recusar ou duvidar de suas entregas.

Também de aplicação para a educação é a questão da atribuição de responsabilidade \cite{coeckelbergh2023}. Se uma pessoa delega à IA a tarefa de produzir um texto, e esse texto é gerado a partir de cálculos probabilísticos sem intencionalidade, de quem é a responsabilidade pelo discurso que ali vai ou pelas eventuais consequências que ele tenha? Supondo que esse texto não seja nem sequer conferido, como muitas vezes ocorre entre estudantes e docentes, como desenvolver a autoria e a criticidade de crianças e jovens diante de um cenário de completa delegação automática da escrita à máquina? Isso pode ser assim formulado:

\begin{quote}
    Para ser responsável, você precisa saber o que está fazendo e acarretando, em retrospecto, saber o que você fez. Ademais, esse assunto possui um aspecto relacional: no caso de seres humanos, esperamos uma explicação de alguém sobre o que fez ou decidiu. Responsabilidade significa, assim, possibilidade de elucidação e explicabilidade \cite[p. 108]{coeckelbergh2023}.
\end{quote}

Não à toa, a identificação de fraudes na produção textual na escola tem sido acompanhada de desconfianças interessantes: textos gramaticalmente corretos demais (além do que o estudante, em certas etapas, pode produzir); e a incapacidade de explicação e elucidação por parte dos ``autores''. \textcite{coeckelbergh2023} aponta que, para alguns tipos de IAG, nem o próprio sistema consegue dar explicações sobre o que fez, nem os humanos por trás dele conseguem explicar o que foi feito pela IA. Em termos de tomada de decisões e de escolhas (aspecto fundamental na abordagem multimodal, por exemplo), itens importantes para a escrita responsável de um texto, o que se tem visto é uma profusão de parágrafos gerados por máquinas e assumidos por humanos, mas sem a devida implicação autoral ou possibilidade de atribuição de responsabilidade. Assim resume \textcite[p. 116, grifos do original]{coeckelbergh2023}: ``Não importando se a IA pode ou não fornecer \textit{diretamente} motivos e explicações, \textit{seres humanos} devem ser capazes de responder à pergunta: `por quê?'".

\section{IA e a criação}
Além de professora e pesquisadora, e antes disso, sou escritora. Nos últimos tempos, ao participar de eventos, uma pergunta direcionada a mim se tornou comum: ``Você usa IA para escrever?''. Há aí, além da curiosidade sobre o que comumente se chama de ``processo criativo'', a desconfiança de que a IA possa potenciar algo que já faço, seja na forma da entrega mais rápida, seja na melhoria da qualidade do texto que publico. Minha resposta, ao menos por enquanto, é ``Não''. Não uso IA para escrever, nem sequer para pesquisar (aliás, função que ela não tem), nem ao menos para revisar um texto que escrevi. O espanto diante da minha resposta tem relação com certo julgamento de que eu esteja perdendo tempo ou dispensando a possibilidade de facilitar meu trabalho.

O uso alienado e indiscriminado, às vezes irresponsável e antiético, da IA tem ensejado debates acalorados e, às vezes, resultado em guias de uso importantes para a calibragem de práticas que ainda não foram totalmente compreendidas, menos ainda suas consequências. Entre as propostas de uso crítico das IAG estão: entregar às inteligências artificiais tarefas menos significativas; não delegar o trabalho criativo a elas; apenas colaborar com elas \cite{coeckelbergh2023}. A preocupação central aí é com o sentido mesmo da vida, isto é, à preservação das atividades muito significativas para as pessoas. A questão pode ser elaborada da seguinte maneira: ``Quais tarefas queremos ou precisamos reservar para os seres humanos e quais poderiam ser os papéis da IA, se houver algum, para nos apoiar nessas tarefas de maneiras eticamente boas e socialmente aceitáveis'' \cite[p. 131]{coeckelbergh2023}. Essa é uma abordagem centrada no humano que me ajuda a responder aos que me perguntam sobre se uso IA para escrever meus contos, poemas e artigos.

Em algumas ocasiões, diante desse questionamento tão atual, me vali da analogia com outras artes e outras produções criativas que parecem ser melhor entendidas como atividades muito significativas para os seres humanos. Perguntei à plateia, em um evento, se alguém ali tocava um instrumento musical, ou se fazia esportes regularmente, ou se preferia dançar. Muitos e muitas levantaram as mãos e disseram sim, com orgulho. Daí eu disse que gosto de escrever e perguntei se alguma dessas pessoas pediria à IA para dançar por elas. Não, ninguém quer isso. Dançar é um prazer, tal como a escrita pode ser, para parte das pessoas. Nesse sentido, a automatização de certos processos na geração do texto pode afetar, justamente, o ponto fundamentalmente humano que separa escrever de gerar textos. \textcite{coeckelbergh2023} menciona ainda a acumulação de riscos tecnológicos, o crescimento das vulnerabilidades humanas e o uso ignorante e imprudente das IA como grandes problemas a serem evitados ou solucionados com urgência, e conclama à construção de uma ética da tecnologia. Isso deve incluir, penso, a distinção cuidadosa entre o que se aprende e se assume no processo de escrever contra o que se delega à máquina para eximir alguém não apenas de certo esforço, mas também da possibilidade da autoria e até mesmo do prazer.

A autoria, aliás, é um dos temas sempre postos sob escrutínio, há décadas. Há textos seminais, nos quais não me deterei, sobre a função autor, a morte do autor e outras questões que levam a pensar na formação de um ser humano capaz de se responsabilizar por um texto escrito, ser moralmente associado a ele e mesmo ser recompensado economicamente. A história da noção de autoria passa por condições históricas e sociais conhecidas e pode-se dizer que jamais esteve pacificada, e não raro porque a chegada de uma nova tecnologia de ler e escrever perturba esse terreno sempre instável. As IAG entram nesse rol. Se a autoria volta ao debate como problema, se alguém assume um texto gerado por IA e não avisa sobre isso, ou mesmo que avise, é interessante lembrar que a percepção dessas nuanças tem importâncias diferentes para as diversas áreas do conhecimento e para suas práticas de produção e publicação, por exemplo. Na área de Saúde ou na de Engenharia, é comum que um artigo seja assinado por grupos de pessoas, às vezes bastante numerosos; pode ser perfeitamente aceitável que uma dissertação de mestrado se componha de três artigos já publicados, inclusive em coautoria com orientadores e colegas. Na área de Letras, até hoje, as coautorias são não apenas raras, como também são vistas com maus olhos, em muitos casos. Num texto assinado por dez autores, o que importa quem o escreveu? Como distinguir quanto cada pessoa contribuiu? Que relevância tem cada contribuição? Aquele/aquela que revisou o texto final e eventualmente o assinou desempenhou uma tarefa menor (e nem deveria estar entre os autores)? É aceitável que um orientador some seu nome a um trabalho de seu orientando, sem tê-lo efetivamente escrito? Em certas áreas, isso não é posto em questão. Por que a IAG incomodaria tanto, e mais uma vez, o processo da composição da(s) autoria(s)?

Esse cenário está descrito a fim de lembrar que a produção de textos e suas condições não se assentam sobre um terreno firme e pacificado. No entanto, neste ensaio, o grau de complexidade das preocupações aumenta porque não trato prioritariamente do ensino superior e do mundo dos adultos com escolarização básica concluída, mas de etapas anteriores, especialmente do ensino médio, espaço de disputas (políticas e econômicas) ferrenhas sobre a formação de uma juventude que ainda tomará lugar no mundo. Quem poderá escrever e avaliar a própria escrita (e a dos outros) sem a experiência direta com o processo de escrever?

\section{Alguns casos de sala de aula}
A famosa agência Press Association usa \textit{bots} para escrever notícias locais. Não é a única. Muitas empresas têm usado IA para gerar muitos textos, desde e-mails em resposta a todo tipo de demanda até as matérias que lemos nas páginas dos jornais. Diante desse cenário de incorporação dessa tecnologia a processos profissionais tão tradicionais, é difícil convencer um jovem de que vale a pena se empenhar no processo de escrever.

Minha sala de aula, hoje, é palco de cenas muito interessantes das quais fazem parte usos de IA um tanto desajeitados. E isso vai aumentar, pelo visto. Na graduação, vivemos episódios com textos de todo tipo, mas especialmente com os Trabalhos de Conclusão de Curso (TCC). Há bem pouco tempo, uma estudante submeteu a uma banca um texto supostamente autoral que se valia de catorze citações da tese de uma professora da instituição. Nenhuma dessas citações era verdadeira. A IAG empregada, diante da demanda de encontrar boas citações para o trabalho, inventou os fragmentos e os atribuiu à autora (essa invenção tem sido chamada de ``alucinação''). Ocorre que essa mesma autora fazia parte da banca e pôde apontar o equívoco, afirmando não reconhecer nenhuma daquelas citações supostamente diretas. Questionada sobre a autoria do trabalho, a aluna assumiu ter usado IA, mas ``só um pouco'', isto é, ela delegara à máquina apenas as citações, segundo confessava. O TCC foi recolhido, a banca foi cancelada e a estudante ganhou uma nova chance não apenas de escrever seu trabalho, mas de aprender a fazer um uso mais refinado da IA, desconfiando sempre do que ela gera. Menos do que uma ocorrência criminosa, como queriam considerar alguns, tratou-se de um uso equivocado da tecnologia, ou o que se poderia considerar a expressão de vulnerabilidade da jovem, ainda incapaz de um uso crítico e ético da IA escolhida.

Em outra ocasião, um estudante do ensino médio entregou à professora de Redação um texto digitado em papel A4. A produção modelo Enem pôde ser feita, extraordinariamente, em casa, ambiente que parece aumentar a tentação do uso de um acelerador como a IAG para a produção de textos. No momento exato da entrega, já na sala de aula, o aluno fez espontaneamente a seguinte ressalva: ``Professora, usei IA para escrever, sim? É só para a senhora não achar que eu posso escrever assim''. Desse episódio decorrem diversas questões de grande interesse para quem atua no ensino de língua: como avaliar esse texto ciborgue? O que considerar entre os critérios de qualidade estabelecidos pelo próprio Enem? Se não foi uma pessoa que deu acabamento àquele texto, qual é a medida da participação do humano e da máquina nesse resultado? O que admite ou assume o estudante quando afirma não ``escrever assim''? ``Assim'' como?

O estudante delega parte de sua tarefa à máquina, e sem pesar, mas avisa à professora ao menos duas coisas: que não assume totalmente a produção e que não tem competência para alcançar um resultado como aquele que consta no papel. Em termos de ``evento de letramento'' \cite{kleiman1995}, de que aprendizagens ou desenvolvimentos o jovem abre mão e escapam à docente? O estudante tem consciência disso? O objetivo de entregar um texto ``correto'' supera a oportunidade de cumprir a tarefa sem ajuda artificial? Se a IA passasse a ser usada e assumida, na escola básica, como ``ferramenta'' para o apoio à produção de textos, o que estaríamos deixando de negociar quanto à aprendizagem dos processos da escrita? Sabemos, há décadas, que os processos de escrita são trabalhosos e transformadores em termos cognitivos \cite{flower2016, kellogg2017}. O que deixaremos de desenvolver, se as IAG forem adotadas tão precocemente?

\section{IA e ensino de língua: questão de finalidade}
Sugiro pôr em jogo aqui dois pontos: o objetivo da aula de redação na educação básica; e a metáfora muito usual da ``ferramenta'' para designar, e mesmo defender ou justificar, o uso de uma IAG como apoio à produção textual.

Hoje em dia, com relativa frequência, escuto que a IAG pode ser uma ferramenta útil para o refinamento da produção escrita, quando não uma aceleradora de textos que levariam horas ou dias para serem escritos. Troca-se a atividade de escrever o texto pela de pensar e produzir o \textit{prompt}. Quanto mais preciso e contextualizado ele for, mais chances de o texto gerado pela IA ser razoável, não raro melhor do que o texto de um humano, mesmo adulto e escolarizado. Já há casos de teses ou outras monografias acadêmicas sendo geradas por IA, com resultados aprovados por bancas, embora ainda desconhecidos em termos de responsabilização e explicabilidade, como dito antes.

Pois bem, é fundamental aqui pensar nos objetivos das disciplinas escolares, a fim de especular sobre o impacto de uma IAG para a produção de textos entregues como ``autorais'' por pessoas que, efetivamente, não os escreveram (ou talvez apenas tenham produzido um bom \textit{prompt}). Se imagino que o objetivo de um professor de História seja ensinar sobre a abolição da escravidão no Brasil; ou que uma professora de Geografia precise envolver seus alunos e ensinar sobre os climas das regiões do país; ou a colega de Biologia mostre à turma as partes que compõem uma célula; talvez seja possível imaginar que a produção textual, seus processos e suas questões não sejam o objetivo desses docentes, até porque a textualidade e a textualização não são seus objetos explícitos de ensino. De maneira lateral, sim, pode-se argumentar que o texto deveria ser uma preocupação de todo professor ou professora, de qualquer área/disciplina, mas o que se sabe, e o que acontece de fato, é que cabe à professora e ao professor de Português (e outras denominações, como Literatura e Redação) o ensino da metalinguagem relacionada à língua, incluindo aspectos da interpretação de textos lidos ou ouvidos, elementos da textualidade, gramática normativa e outros itens ditos como ``conteúdos'' da matéria. Há, sim, um equívoco técnico em considerar língua por escrita, já que a escrita é um código convencionado que pouco tem a ver com a língua falada, por exemplo, mas a disciplina escolar conhecida como Português frequentemente aborda questões de leitura e escrita como ``língua'', e reforça a ideia de que escrever mal se relaciona com não saber português. Minha especulação trata, então, dessa distorção operante nas escolas, tomando-a como um fato.

Se a professora de Biologia pode, eventualmente, se preocupar menos com a efetiva autoria de um trabalho sobre células; se o professor de História pode, igualmente, dar menos atenção ao texto sobre a abolição da escravidão, avaliando mais o conteúdo do que os processos de produção textual e seus resultados, o mesmo não se pode dizer da professora de Português, a quem cabe avaliar justamente o processo de produção textual e, claro, o texto resultante desse processo, que, inclusive, admite a refação como oportunidade de aprendizagem. Isto é, um texto entregue em primeira versão, avaliado pela professora, devolvido com comentários, entregue em segunda versão etc. é exatamente o objeto da aula de língua, cujo objetivo toca justamente orientar e ensinar sobre mecanismos de produção escrita, o que passa pela entrega de uma versão eivada de imperfeições ou ``errada'', se for o caso, em relação a parâmetros como a variedade culta da língua. É no ``erro'' que a professora flagra uma dúvida, um ponto a ser discutido, uma possibilidade de avanço. Sem errar, o estudante camufla ou escamoteia um episódio importante do processo de aprender a escrever textos melhores, assim como ficam silenciadas as chances de debate sobre o que e como fazer para produzir novas versões (e melhoradas).

Por que as pessoas fazem isso? Por que elas têm o desejo ou cedem à tentação de driblar o processo de aprendizagem, usando ``ferramentas'' que calculam melhor do que elas? Se a escola é justamente o espaço do desenvolvimento de conhecimentos importantes para a plena cidadania, com e sem uso de tecnologias digitais \cite{soares2002}, por que razão alguém deveria forjar um texto gerado por IA e entregá-lo como se fosse seu?\footnote{Esse é um aspecto que também tem chamado a atenção nas escolas, mas no qual não me deterei aqui. A noção de plágio vem sendo empregada nesse contexto, só que de maneira equivocada. A fraude na autoria do texto gerado por IA tem mais a ver com falsidade ideológica do que com plágio. O texto gerado não assume autoria, embora opere calculando prováveis resultados mais comuns ou mais frequentes, com base em modelos linguísticos geralmente satisfatórios. Assumir um texto gerado por máquina sem o trabalho corporal e intelectual de escrevê-lo é uma espécie de fraude que tem a ver com enganar o outro e enganar-se a si. Passaremos a avaliar então apenas a autoria dos \textit{prompts}?} Se uma disciplina escolar tem por objetivo principal fazer com que crianças e jovens desenvolvam suas escritas, isto é, o conteúdo é também a forma e o exercício, que impacto a IAG pode ter nesse percurso que contorna o episódio da aprendizagem? Pode-se dizer que a IA, no caso da aula de Português, atua exatamente no intervalo entre a tentativa do aluno/da aluna e a aprendizagem que pode advir do retorno ou da negociação a partir das versões do texto. O que é da aula de Redação sem o texto tentativo?

Certamente haverá quem argumente que a IAG pode ser usada no refinamento e na revisão dos textos, mesmo os de jovens aprendizes da escola básica. Ou que, ao gerar um texto, a IAG pode ensinar por meio do \textit{feedback}. A questão aí recai sobre outro elemento do ensino e da aprendizagem da escrita: quem produzirá o \textit{prompt} que orienta a revisão, sem saber o que revisar? Se a IAG é verbocentrada, o que é necessário aprender antes de aprender a escrever \textit{prompts}? Diante de um resultado possivelmente satisfatório gerado pela IA (um texto escrito ou ``apenas'' revisado), quem o avaliará ou decidirá sobre esse resultado e sua qualidade? Quem terá condições de solicitar ajustes e uma nova versão? Diante desse impasse (ou dilema?), em que interessa um texto perfeito (ou eventualmente perfeito), se a dúvida e o erro é que geralmente indiciam o aprendido e o ainda por aprender? Vamos desistir de aprender?

O que quero destacar neste ensaio é que a IAG atinge em cheio o ``conteúdo'' mesmo da aula de Português/Redação. Ela pode gerar como resultado de um \textit{prompt} um texto bastante bom, mas sem indícios que ajudariam no processo de aprendizagem de uma pessoa que não sabe, que duvida, que oscila, que vacila, que ainda não está certa de algumas regras ou especificidades. E as crianças e os jovens frequentemente são essas pessoas, tecnologicamente vulneráveis. A IAG, portanto, não tem, na aula de Redação, um papel de ferramenta, como algo que aumenta a força, o alcance ou auxilia numa tarefa desafiadora, por vezes impossível de executar. A IAG, se usada nesse contexto, ocupa o espaço do objeto mesmo da aula, tornando ocluso o item que deveria servir de ponte para a aprendizagem. Em outras palavras: se a IAG gera um texto com zero problema de pontuação, assumido e entregue por uma criança de 12 anos, temos um adolescente sem oportunidade de negociação sobre vírgulas e pontos, coesão ou argumentação, e muito provavelmente sem a chance de desenvolver o conhecimento necessário para avaliar se o texto da IAG está adequado nesses aspectos e em todos os outros.

\section{Considerações finais ou Não sou professora de \textit{prompt}}
Já se pode encontrar um tipo de IA desenvolvida para reconhecer o uso de IA num texto. Há instituições escolares que, apavoradas com a profusão de textos fraudados ou agora frequentemente desconfiadas de todas as entregas dos alunos e das alunas, começam a sugerir ou a compulsar que docentes fiscalizem a autoria dos textos, obrigando-os a usar ferramentas de identificação de IAG, tornando-os caçadores de algo que nem sequer foi devidamente discutido (e, claro, a se distanciarem, cada vez mais, de seu objetivo precípuo). Grandes universidades começam a publicar seus guias com orientações de usos éticos e adequados de IAG \cite{cnpq_integridade_2026}, materiais que deveriam servir de base para o debate na escola básica, sempre a considerar, claro, que um contexto de ensino e aprendizagem com crianças e jovens tem suas importantes particularidades. Há instituições escolares que propõem aulas de \textit{prompts}, a alunos e professores, alegando uma jornada rumo ao futuro, sem ao menos considerar que a discussão passa muito além do comando.

Em \textcite{dejesus2025}, mostramos resultados de redações de estudantes do ensino fundamental, ponderando que, de fato, uma IAG pode gerar textos adequados, mas, como sabemos, sem intencionalidade ou moralidade alguma. \textcite{buzato2024} mostram como a redação do Enem e sua fórmula podem ser afetadas pela IAG, tão distante ela está de um polo criativo das possibilidades humanas. Já há editoras de livros e de revistas que apresentam suas publicações sob selos de ``Livre de IA'' \textit{(AI free)} ou ``Escrito por seres humanos'', como uma espécie de garantia de qualidade ou, ao menos, de autoria orgânica. Aliás, já aparecem memes e charges rotulando escritores como ``artesanais'' e outros adjetivos ligados a processos não intermediados ou realizados por inteligências artificiais.

O que eu quis considerar aqui, sem qualquer uso de IAG, sendo uma veterana pesquisadora e usuária das tecnologias digitais, foi o erro como possibilidade de diagnóstico, como sempre pôde ser (embora, na escola básica, ele seja frequentemente punido e raramente seja considerado como travessia); ou o erro como o novo índice de humanidade na produção textual; a autoria em disputa, como sempre esteve, mas agora entre humanos e máquinas, omissões e fantasmagorias; a mediação algorítmica desnecessária, e mesmo prejudicial, em alguns contextos de ensino e aprendizagem; a necessidade de que se tenha claramente em mente o objetivo precípuo da aula de Redação e suas dinâmicas que dependem de versões de textos; ou, quem sabe, que essas finalidades sejam revistas, mas não para o atrofiamento de nossas competências e de nosso poder semiótico \cite{kress2003}; a consideração acolhedora das versões textuais como esforços que muitas pessoas quererão evitar ou saltar; o uso 0\% de IAG como um marco entre as possibilidades atuais de processos de escrita, inclusive e principalmente para profissionais capazes da avaliação de textos e da autoavaliação de seus textos; a consciência de que, mesmo que uma pessoa não precise escrever muito ao longo da vida, precisará ao menos ler e avaliar outras escritas, talvez ou eventualmente até mesmo os textos gerados por IA, e deverá, para fazer isso bem, ter aprendido a escrever ou sobre a escrita; a existência de pessoas, de todas as idades, que desejam aprender a escrever ou escrever, e que prefiram e escolham não delegar a tarefa à máquina, ou seja, é não abrir mão de um prazer e de um sentido vital; a necessidade irrefutável de entender as IAG, aprender sobre elas, saber usá-las, compreender sua interação conosco e sua interferência na vida social e na educação, nem que seja para decidir por, em certos contextos, não usá-las ou, como em certos artefatos do consumo, usá-las com moderação, sem deixar escapar o que é preciso aprender e ensinar.

\section*{Agradecimentos}
Agradecimentos à colega Jenny Teresita Guerra González, da Universidade Nacional Autônoma do México (UNAM), pelo diálogo sempre de vanguarda sobre tecnologia e livro.




\printbibliography\label{sec-bib}
% if the text is not in Portuguese, it might be necessary to use the code below instead to print the correct ABNT abbreviations [s.n.], [s.l.]
%\begin{portuguese}
%\printbibliography[title={Bibliography}]
%\end{portuguese}



\begin{dataavailability}
\txtdataavailability{nodata} % options: dataavailable, dataonly, databody, datanotav, nodata
\end{dataavailability}




\end{document}

